
\section{The \name~Benchmark}
\label{sec:methods}

In this section, we detail the design philosophy and construction of the evaluation dataset.
The dataset is designed to provide a broad set of basic evaluations for reward models, covering chat, instruction following, coding, safety, and other important metrics for fine-tuned language models.
The \name~dataset contains a combination of existing evaluation prompt-completion pairs, and those curated for this project.

A good reward function, and therefore a good RM broadly, is one that stably assigns credit to the classes of good or bad content.%\footnote{There are more considerations on how to \textit{use} a RM, but the initial notion of quality should be one that agrees with curated data. Next, we can evaluate which RMs are best for downstream tasks such as RLHF.}
Given one verified answer that is better than another for factual or clear qualitative reasons (e.g. typos), a good reward model will choose the correct one 100\% of the time.
To evaluate this, each datapoint consists of a \texttt{prompt} and two completions, \texttt{chosen} and \texttt{rejected}.
For each prompt, the score of the reward model is computed. The prompt is then categorized as a win if the score of the prompt with the verified chosen completion is higher than that of the verified rejected completion, as shown in Fig.~\ref{fig:method}.
Finally, we report accuracy for each subset as the percentage of wins. 
For all the section scores of~\name~(e.g. \texttt{Chat} or \texttt{Safety}) except \texttt{Prior Sets}, the average score is weighted per-prompt in the requisite subsets.

\subsection{\name~Dataset}
\label{sec:subsets}

The benchmark is broken down into five sections from different subsets -- the first four compose the~\name~dataset described in this section.
We have broken down the dataset into these subsections to create one final~\name~score in order to reasonably weigh different aspects of an RM's performance.
The RewardBench dataset is released under the ODC-BY license\footnote{ODC-BY: \url{https://opendatacommons.org/licenses/by/1-0/}} and the code is released under Apache 2.0\footnote{Apache 2.0: \url{https://www.apache.org/licenses/LICENSE-2.0}}.
The summary of the dataset is shown in Tab.~\ref{table:dataset_char} (see appendix~\ref{appdx:data_details} for full details) At a high level, the subsets consist of the following:

\begin{enumerate}[leftmargin=0.5cm]
    \item \textbf{Chat}: Testing a reward model's basic ability to distinguish a thorough and correct chat response in open-ended generation.
    Prompts and chosen, rejected pairs are selected from AlpacaEval~\citep{alpaca_eval} and MT Bench~\citep{zheng2023judging}, two popular open-ended chat evaluation tools. 
    \item \textbf{Chat Hard}: Testing a reward model's abilities to understand trick questions and subtly different instruction responses. 
    Prompts and chosen, rejected pairs are selected from MT Bench examples with similar ratings and adversarial data specifically for fooling LLM-as-a-judge tools from LLMBar's evaluation set~\citep{zeng2023llmbar} (reformatted for RMs).
    \item \textbf{Safety}: Testing the models' tendencies to refuse dangerous content and to avoid incorrect refusals to similar trigger words. 
    Prompts and chosen, rejected pairs are selected from custom versions of the datasets XSTest~\citep{rottger2023xstest}, Do-Not-Answer~\citep{wang2023donotanswer}, and examples from an in-development refusals dataset at AI2, where the chosen response is a refusal and the rejected is harmful text of either dangerous or offensive nature.
    \item \textbf{Reasoning}: Evaluating the models code and reasoning abilities. 
    Code prompts are created by reformatting HumanEvalPack examples with correct code as chosen and rejected as one with bugs~\citep{muennighoff2023octopack}.
    Reasoning prompts pair reference answers with incorrect model generations from the PRM800k dataset~\citep{lightman2023let}.
    
    \begin{table}[t]
\caption{
Top-20 open models on \name. 
Evaluating many RMs shows that there is still large variance in RM training and potential for future improvement across the more challenging instruction and reasoning tasks.
Icons refer to model types: Sequence Classifier (\sequenceclf), Direct Preference Optimization (\dpo), Custom Classifier (\customclf), Generative Model (\generative), and a random model (\random).
}
\vspace{0.5em} 
{\footnotesize
\centering
\begin{tabular}{lrrrrrr}
%\begin{tabular}{lcccccc}
%\toprule
% Reward Model & \rot{Average} & \rot{Chat} & \rot{Chat Hard} & \rot{Safety} & \rot{Reasoning} & \rot{Prior Sets} \\
Reward Model & \thead{Score} & \thead{Chat} & \thead{Chat\\Hard} & \thead{Safety} & \thead{Reason} & \thead{Prior\\Sets} \vspace{-3pt} \\ 
% Reward Model & \multicolumn{2}{l}{Average}      & \multicolumn{2}{l}{Chat Hard}        & \multicolumn{2}{l}{Reasoning}            \\
%             & \multicolumn{2}{r}{Chat}          &         \multicolumn{2}{r}{Safety}     &    \multicolumn{2}{r}{Prior Sets}        \\
\midrule
\href{https://huggingface.co/RLHFlow/ArmoRM-Llama3-8B-v0.1}{\customclf RLHFlow/ArmoRM-Llama3-8B-v0.1} & \textbf{89.0} & 96.9 & \textbf{76.8} & \textbf{92.2} & \textbf{97.3} & 74.3 \\
% \href{https://huggingface.co/google/gemini-1.5-pro-0514}{\generative google/gemini-1.5-pro-0514} & 88.1 & 92.3 & 80.6 & 87.5 & 92.0 & - \\
\href{https://huggingface.co/RLHFlow/pair-preference-model-LLaMA3-8B}{\customclf RLHFlow/pair-preference-model-LLaMA3-8B} & 85.7 & 98.3 & 65.8 & 89.7 & 94.7 & 74.6 \\
% \href{https://huggingface.co/openai}{\generative openai/gpt-4-0125-preview} & 84.3 & 95.3 & 74.3 & 87.2 & 86.9 & 70.9 \\
% \href{https://huggingface.co/openai}{\generative openai/gpt-4-turbo-2024-04-09} & 83.9 & 95.3 & 75.4 & 87.1 & 82.7 & 73.6 \\
\href{https://huggingface.co/sfairXC/FsfairX-LLaMA3-RM-v0.1}{\sequenceclf sfairXC/FsfairX-LLaMA3-RM-v0.1} & 83.6 & \textbf{99.4} & 65.1 & 87.8 & 86.4 & 74.9 \\
% \href{https://huggingface.co/openai}{\generative openai/gpt-4o-2024-05-13} & 83.3 & 96.6 & 70.4 & 86.7 & 84.9 & 72.6 \\
\href{https://huggingface.co/openbmb/Eurus-RM-7b}{\sequenceclf openbmb/Eurus-RM-7b} & 81.6 & 98.0 & 65.6 & 81.2 & 86.3 & 71.7 \\
\href{https://huggingface.co/Nexusflow/Starling-RM-34B}{\sequenceclf Nexusflow/Starling-RM-34B} & 81.4 & 96.9 & 57.2 & 88.2 & 88.5 & 71.4 \\
% \href{https://huggingface.co/Anthropic}{\generative Anthropic/claude-3-opus-20240229} & 80.7 & 94.7 & 60.3 & 89.1 & 78.7 & - \\
\href{https://huggingface.co/weqweasdas/RM-Mistral-7B}{\sequenceclf weqweasdas/RM-Mistral-7B} & 79.3 & 96.9 & 58.1 & 87.1 & 77.0 & \textbf{75.3} \\
\href{https://huggingface.co/hendrydong/Mistral-RM-for-RAFT-GSHF-v0}{\sequenceclf hendrydong/Mistral-RM-for-RAFT-GSHF-v0} & 78.7 & 98.3 & 57.9 & 86.3 & 74.3 & 75.1 \\
\href{https://huggingface.co/stabilityai/stablelm-2-12b-chat}{\dpo stabilityai/stablelm-2-12b-chat} & 77.4 & 96.6 & 55.5 & 82.6 & 89.4 & 48.4 \\
\href{https://huggingface.co/Ray2333/reward-model-Mistral-7B-instruct-Unified-Feedback}{\sequenceclf Ray2333/reward-model-Mistral-7B-instruct...} & 76.9 & 97.8 & 50.7 & 86.7 & 73.9 & 74.3 \\
\href{https://huggingface.co/allenai/tulu-2-dpo-70b}{\dpo allenai/tulu-2-dpo-70b} & 76.1 & 97.5 & 60.5 & 83.9 & 74.1 & 52.8 \\
% \href{https://huggingface.co/PoLL/gpt-3.5-turbo-0125_claude-3-sonnet-20240229_meta-llama/Llama-3-70b-chat-hf}{\generative PoLL/gpt-3.5-turbo-0125\_claude-3-sonnet-20240229...} & 75.6 & 95.3 & 54.1 & 79.5 & 73.5 & - \\
\href{https://huggingface.co/meta-llama/Meta-Llama-3-70B-Instruct}{\generative meta-llama/Meta-Llama-3-70B-Instruct} & 75.4 & 97.6 & 58.9 & 69.2 & 78.5 & 70.4 \\
\href{https://huggingface.co/prometheus-eval/prometheus-8x7b-v2.0}{\generative prometheus-eval/prometheus-8x7b-v2.0} & 75.3 & 93.0 & 47.1 & 83.5 & 77.4 & - \\
% \href{https://huggingface.co/Anthropic}{\generative Anthropic/claude-3-sonnet-20240229} & 75.0 & 93.4 & 56.6 & 83.7 & 69.1 & 69.6 \\
\href{https://huggingface.co/NousResearch/Nous-Hermes-2-Mistral-7B-DPO}{\dpo NousResearch/Nous-Hermes-2-Mistral-7B-DPO} & 74.8 & 92.2 & 60.5 & 82.3 & 73.8 & 55.5 \\
\href{https://huggingface.co/mistralai/Mixtral-8x7B-Instruct-v0.1}{\dpo mistralai/Mixtral-8x7B-Instruct-v0.1} & 74.7 & 95.0 & 64.0 & 73.4 & 78.7 & 50.3 \\
\href{https://huggingface.co/upstage/SOLAR-10.7B-Instruct-v1.0}{\dpo upstage/SOLAR-10.7B-Instruct-v1.0} & 74.0 & 81.6 & 68.6 & 85.5 & 72.5 & 49.5 \\
% \href{https://huggingface.co/Anthropic}{\generative Anthropic/claude-3-haiku-20240307} & 73.5 & 92.7 & 52.0 & 82.1 & 70.6 & 66.3 \\
\href{https://huggingface.co/HuggingFaceH4/zephyr-7b-alpha}{\dpo HuggingFaceH4/zephyr-7b-alpha} & 73.4 & 91.6 & 62.5 & 74.3 & 75.1 & 53.5 \\
\href{https://huggingface.co/allenai/tulu-2-dpo-13b}{\dpo allenai/tulu-2-dpo-13b} & 73.4 & 95.8 & 58.3 & 78.2 & 73.2 & 49.5 \\
\href{https://huggingface.co/0-hero/Matter-0.1-7B-boost-DPO-preview}{\dpo 0-hero/Matter-0.1-7B-boost-DPO-preview} & 73.4 & 91.1 & 61.0 & 66.3 & 83.9 & 55.7 \\
\href{https://huggingface.co/prometheus-eval/prometheus-7b-v2.0}{\generative prometheus-eval/prometheus-7b-v2.0} & 72.4 & 85.5 & 49.1 & 78.7 & 76.5 & - \\
\href{https://huggingface.co/HuggingFaceH4/starchat2-15b-v0.1}{\dpo HuggingFaceH4/starchat2-15b-v0.1} & 72.1 & 93.9 & 55.5 & 65.8 & 81.6 & 55.2 \\
% \href{https://huggingface.co/HuggingFaceH4/zephyr-7b-beta}{\dpo HuggingFaceH4/zephyr-7b-beta} & 71.8 & 95.3 & 62.7 & 61.0 & 77.9 & 52.2 \\
% \href{https://huggingface.co/allenai/tulu-2-dpo-7b}{\dpo allenai/tulu-2-dpo-7b} & 71.7 & 97.5 & 56.1 & 73.3 & 71.8 & 47.7 \\
% \href{https://huggingface.co/jondurbin/bagel-dpo-34b-v0.5}{\dpo jondurbin/bagel-dpo-34b-v0.5} & 71.5 & 93.9 & 55.0 & 61.5 & 88.9 & 44.9 \\
% \random \textit{Random} & 50.0 & 50.0 & 50.0 & 50.0 & 50.0 & 50.0 \\
\bottomrule
\end{tabular}}
% \vspace{0.5em} 
% \caption{
% Top-20 open models on \name. 
% Evaluating many RMs shows that there is still large variance in RM training and potential for future improvement across the more challenging instruction and reasoning tasks.
% Icons refer to model types: Sequence Classifier (\sequenceclf), Direct Preference Optimization (\dpo), Custom Classifier (\customclf), Generative Model (\generative), and a random model (\random).
% }
\label{table:top20_results}
\end{table}

    \item \textbf{Prior Sets}\footnote{For the final RewardBench score, we weigh the Prior Sets category at 0.5 weight of the others due to multiple factors: noise, lack of clearly defined tasks, etc. The dataset is found here: \url{https://huggingface.co/datasets/allenai/preference-test-sets} }: For consistency with recent work on training reward models, we average performance over test sets from existing preference datasets.
    We use the Anthropic Helpful split~\citep{bai2022training} (the only multi-turn data), the Anthropic HHH subset of BIG-Bench~\citep{askell2021general}, a curated subset of the test set from the Stanford Human Preferences (SHP) Dataset~\citep{pmlr-v162-ethayarajh22a}, and OpenAI's Learning to Summarize Dataset~\citep{stiennon2020learning}.
\end{enumerate}




\subsection{\name~Scoring}

\name~is scored via accuracy.
For each prompt-chosen-rejected trio, we infer the score the RM assigns for the prompt-chosen and prompt-rejected pairs
%\footnote{For some reward models, such as PairRM and SteamSHP, their intended use is with pairwise inputs, so we evaluate in that manner following the original source code.} 
then assign a true classification label when the chosen score is higher than rejected, as highlighted in Fig.~\ref{fig:method}.
Details on computing scores for classifiers and DPO models is in Sec.~\ref{sec:back}.
Given the binary classification task, a random model achieves a result of 50\%.
%On many subsets, models achieve at or well below the random baseline, indicating substantial areas of progress in reward models. \vp{Todo: @ Nathan rephrase}
In order to create a representative, single evaluation score, we perform a mixture of averaging across results.
For the sections detailed in Sec.~\ref{sec:subsets} except for Reasoning, we perform per-prompt weighted averaging across the subsets to get the normalized section scores.
For example, in Chat we take a weighted average of the AlpacaEval and MT Bench sets based on the number of prompts.
For Reasoning, we increase the weight of the PRM-Math subset so code and math abilities are weighed equally in the final number.
For Prior Sets, we take an unweighted average over the subsets due to the large disparity in dataset sizes.
Once all subsets weighted averages are achieved, the final~\name~score is the weighted average across the section scores (Prior Sets at 0.5 weight).